%
% File: discussion.tex
%
\chapter{Discussion and Conclusion}
\label{chap:discussion}
\setlength{\parskip}{1em}
%
\section{Theoretical Implications}

The synthesis presented in Chapter~\ref{chap:findings} has three implications for multimedia learning theory. First, it supports the argument by \citet{twabu2025} that CTML and CLT, while foundational, require extension to accommodate the adaptive, learner-responsive behaviour of generative AI. The classical principles remain valid as design heuristics, but they were articulated for static or pre-authored multimedia, not for systems that produce content in real time in response to learner input. The frameworks need conceptual room for AI as an active participant in the learning environment, not merely as a delivery channel.

Second, the review confirms the centrality of the active-processing assumption. The most reliable finding across the corpus --- supported by the experimental work of \citet{lawson2024}, by the interview data of \citet{buhler2024}, and by the design choices in the tool of \citet{osman2024} --- is that visual content alone is insufficient. The cognitive activity surrounding the visual is what produces learning. This finding sets a clear constraint on AI-generated visual explanations: they must be designed to provoke cognitive engagement, not to substitute for it.

Third, the review identifies adaptive cognitive load management as a productive frontier for theoretical work. AI systems can in principle measure and respond to learner load in ways that conventional multimedia cannot. Whether they should, and under what conditions, is a question that current theory does not yet answer. This is the most promising direction for theoretical extension.

\section{Practical Implications}

For educators and instructional designers working with learners with low language proficiency or learning disabilities, the review yields several practical recommendations. First, the choice of AI tool should be guided by its alignment with multimedia learning principles, not by its technological novelty. A tool that respects the modality, signalling, and coherence principles will, on the available evidence, support learning more reliably than a more powerful tool that violates them. Second, AI-generated visuals should be embedded in active tasks: prompts that ask the learner to summarise, predict, compare, or critique the AI output, rather than simply receive it. The tasks need not be elaborate, but they must be present.

Third, learners should be supported in developing critical engagement with AI outputs. The accuracy concerns documented in \citet{buhler2024} and the bias concerns raised by \citet{ahmad2025} are not problems that pedagogy can solve in isolation, but pedagogy can equip learners to detect and respond to them. Fourth, the educator's role remains central. AI-generated visual explanations should be treated as material that the educator selects, contextualises, and validates, not as autonomous instructional agents.

Fifth, institutions adopting these tools at scale should consider equity of access. The benefits documented in the literature accrue to learners with stable connectivity, capable devices, and where applicable, paid subscriptions. Without institutional provision, the technology may widen rather than narrow the access gap for the very learners it claims to support.

\section{Limitations of the Review}

Several limitations of this review should be acknowledged. The search was restricted to English-language publications, which excludes potentially relevant work in Russian, Mandarin, Spanish, and other languages with substantial educational research traditions. The date range of 2023--2025 reflects the post-GPT-4 era of multimodal generative AI but excludes earlier work that might illuminate longer-term trajectories. Source extraction was conducted by the two authors of this review without independent third-party verification, which introduces potential for selection and interpretation bias; future iterations should employ formal inter-rater reliability assessment. The corpus of ten studies is small relative to the breadth of the topic, although deliberately so: the review prioritised depth of engagement over breadth of coverage. The thematic synthesis approach surfaces patterns but does not produce statistical effect-size estimates, which a meta-analysis on a more homogeneous sub-question could provide. Finally, the review treats learners with low language proficiency and learners with learning disabilities as a single population for analytic convenience, but these groups have distinct needs that future work should differentiate more carefully.

\section{Directions for Future Research}

Four research directions follow from the synthesis. First, longitudinal studies are needed. The corpus is dominated by short-term evaluations and design studies; the cumulative effects of sustained exposure to AI-generated visual explanations on learner agency, metacognition, and language development remain unexamined. Second, differential-effects research is needed: the review has treated the target populations together, but evidence on how AI-generated visual explanations function for learners with dyslexia versus ASD versus low language proficiency, and at different ages, is currently sparse. Third, controlled studies of adaptive load management are needed to test the theoretical proposals advanced by \citet{twabu2025} and others. Fourth, culturally responsive design research is needed: most of the corpus reflects Anglophone, Western, well-resourced contexts, and the transferability of findings to other contexts --- including Kazakhstani higher education, where the present review is situated --- cannot be assumed. Cross-cultural comparative studies and studies grounded in non-Western educational traditions would substantially strengthen the field's external validity.

\section{Conclusion}

AI-generated visual explanations are neither a panacea for inclusive education nor a passing trend. They are a technological capability whose pedagogical potential is real, partially evidenced, and conditional on careful design and institutional support. For learners with low language proficiency or learning disabilities, the technology offers a route around the linguistic bottleneck that has long constrained their access to complex content. Realising that potential will require a research base more integrated than the current literature, design practice more anchored in theory than in novelty, and institutional commitment more attentive to equity than to efficiency. The frameworks and evidence reviewed here suggest that this is feasible.
